% Part: lambda-calculus
% Chapter: lambda-definability
% Section: lambda-definable-recursive

\documentclass[../../../include/open-logic-section]{subfiles}

\begin{document}

\olfileid{lam}{dfl}{ldr}
\olsection{Fungsi kang \usetoken{S}{lambda definable} iku Rekursif}

Ora mung kabeh fungsi rekursif parsial iku !!{lambda definable},
kosokbaline uga bener. Tegese, kabeh fungsi kang
!!{lambda definable} iku rekursif parsial.

\begin{thm}
  \ollabel{thm:lambda-computable} Yen fungsi parsial $f$ iku
  !!{lambda definable}, mula fungsi iku rekursif parsial.
\end{thm}

\begin{proof}
  Kita mung menehi garis gedhe buktine. Sepisanan, kita
  aritmetisasi suku-$\lambd$, yaiku menehi wilangan G\"odel
  marang suku-$\lambd$ kanthi sistematis, nganggo cara lumrah
  pangodean rerangkening liwat pangkat wilangan prima. Banjur
  kita netepake fungsi rekursif parsial $\fn{normalize}(t)$ kang
  nampa wilangan G\"odel~$t$ saka suku lambda minangka argumen.
  Fungsi iku mbalekake wilangan G\"odel saka wujud normal
  yen ana, lan ora nduweni nilai yen ora ana. Banjur tetepna
  rong fungsi rekursif parsial $\fn{toChurch}$ lan
  $\fn{fromChurch}$ kang ngowahi wilangan asli dadi wilangan
  G\"odel saka numeral Church kang cocog, lan kosokbaline.

  Kanthi fungsi-fungsi rekursif iki, kita bisa netepake fungsi~$f$
  minangka fungsi rekursif parsial. Ana suku-$\lambd$~$F$ kang
  !!{lambda define}s~$f$. Kanggo ngitung $f(n_1, \dots, n_k)$,
  dhisik goleka wilangan G\"odel saka numeral Church kang cocog
  nganggo $\fn{toChurch}(n_i)$, banjur sambungna kode-kode iki
  marang $\Gn{F}$ kanggo entuk wilangan G\"odel saka suku
  $F \num{n_1}\dots\num{n_k}$. Saiki gunakna $\fn{normalize}$
  marang wilangan G\"odel iki. Yen $f(n_1, \dots, n_k)$ nduweni
  nilai, $F \num{n_1}\dots\num{n_k}$ nduweni wujud normal
  (kang kudu numeral Church); yen ora, suku iki ora nduweni
  wujud normal (saéngga
  \[\fn{normalize}(\Gn{F\num{n_1}\dots\num{n_k}})\]
  uga ora nduweni nilai). Pungkasan, gunakna $\fn{fromChurch}$
  marang wilangan G\"odel saka suku kang wis dinormalake.
\end{proof}

\end{document}
